\documentclass[a4paper,12pt]{article}
\usepackage[utf8]{inputenc}
\usepackage{graphicx}
\usepackage{xcolor}
\usepackage{enumitem}
\usepackage{caption}
\usepackage[most]{tcolorbox}
\usepackage{float}
\usepackage[a4paper,
  top=3.1cm,
  bottom=3.2cm,
  left=3cm,
  right=3cm
]{geometry}
\directlua{dofile("catalogue.lua")}

\newlength{\cataloguewidth}
\setlength{\cataloguewidth}{\textwidth}

\newcommand{\ruleexamples}[2]{
\begin{center}
\begin{tabular}{cc}
\includegraphics[width=0.35\textwidth]{#1} &
\includegraphics[width=0.35\textwidth]{#2}
\end{tabular}
\end{center}
}

\newtcolorbox{excludedrule}{
    colback=black!5,
    colframe=black!20,
    coltext=black!45,
    boxrule=0.5pt,
    arc=1mm,
    left=6pt,
    right=6pt,
    top=5pt,
    bottom=5pt,
    before skip=3pt,
    after skip=3pt
}


\newcommand{\rulemeta}[3]{
\smallskip\\
\textcolor{orange}{Event:} #1 \quad
\textcolor{orange}{Condition:} #2 \quad
\textcolor{orange}{Stimulus:} #3
}


\title{Experiment Design and Rule Catalog}
\author{}
\date{}

\begin{document}


\vspace*{-2cm}
{\let\newpage\relax\maketitle}

\vspace*{-1cm}
\section*{Design}

The experiment consists of six sessions with five blocks each, one per rule family: Arithmetic, Attraction, Recoloring, Expansion, and Occlusion.

Each block contains one rule family and one of two contexts (Figure \ref{fig:task-contexts}):

\begin{itemize}
    \item \textbf{Previous rule (Inference):} compare the current rule with the rule from the previous trial.
    \item \textbf{First/Memorized rule (Application):} compare the current rule with the rule introduced at the start of the block.
\end{itemize}

Each trial presents two examples of the same hidden rule. A block contains one initial rule trial followed by eight same/different decisions, balanced to four \textit{same} and four \textit{different} trials with a maximum run length of three identical responses.

Contexts alternate within sessions and the starting context alternates across sessions. Rule families are assigned such that each family occurs equally often in both contexts across the six sessions.

Frame color is counterbalanced across participants: Variant A (odd participant numbers) uses yellow for Inference and cyan for Application; Variant B (even participant numbers) reverses this mapping.

Trials last 10\,s and each block is preceded by 10\,s fixation. Including the final fixation and 10\,s of dummy scans (e.g. TR = 2\,s), each session requires 520\,s (8\,min 40\,s; 260 volumes).

\newpage
\newgeometry{top=3cm,bottom=3cm,left=1.2cm,right=1.2cm}
\thispagestyle{empty}

\begin{figure}[H]
\centering

{\large\bfseries Application / Memorized rule}

\medskip

\setlength{\tabcolsep}{2pt}
\begin{tabular}{cc}
\includegraphics[width=0.49\textwidth]{images/memorize_first_rule.png} &
\includegraphics[width=0.49\textwidth]{images/memorize_different.png} \\
\small Initial trial &
\small Decision trial
\end{tabular}

\vfill

\vspace{1cm}

{\large\bfseries Inference / Previous rule}

\medskip

\begin{tabular}{cc}
\includegraphics[width=0.49\textwidth]{images/first_rule.png} &
\includegraphics[width=0.49\textwidth]{images/previous_same.png} \\
\small Initial trial &
\small Decision trial
\end{tabular}

\vspace{0.2cm}

\caption{Example screens for the two experimental contexts (Variant A).}
\label{fig:task-contexts}
\end{figure}

\newpage
\restoregeometry


\directlua{render_catalogue("images/rules")}


\end{document}